\begin{proposition}[6.15.6]
\label{IV.2.s6.15.6-fr}
Sean $k$ un cuerpo y $X$ un $k$-preesquema. Si $X$ es normal,
entonces, para toda extensión $k'$ de $k$,
$X'=X\otimes_k k'$ es geométricamente unirramoso.
\end{proposition}

\begin{proof}
Se sabe que $k'$ es una extensión algebraica de una extensión
puramente trascendente $k_0$ de $k$ y que, si $k''$ es la mayor
extensión separable de $k_0$ contenida en $k'$, $k'$ es una
extensión radicial de $k''$ y $k''$ una extensión separable de
$k$. Se sabe (6.14.2) que $X''=X\otimes_k k''$ es normal; como
$X\otimes_k k'=X''\otimes_{k''}k'$, se ve que es posible limitarse al caso
en que la extensión $k'$ de $k$ es \emph{radicial}. Además (\textbf I, 3.6.5),
se puede suponer que $X=\operatorname{Spec}(A)$, donde $A$ es un
anillo \emph{local íntegro e íntegramente cerrado} (puesto que $X$
es normal). El morfismo de proyección $f:X'\to X$ es un
homeomorfismo, ya que
$\operatorname{Spec}(k')\to\operatorname{Spec}(k)$ es un
homeomorfismo universal (2.4.5); como
$X'=\operatorname{Spec}(A')$, donde $A'=A\otimes_k k'$, se ve
por tanto que $A'$ es un anillo local cuyo nilradical
$\mathfrak N'$ es el único ideal primo minimal. De ahí que
$X'_{\mathrm{red}}=\operatorname{Spec}(A'_0)$, donde
$A'_0=A'/\mathfrak N'$ es un anillo local íntegro. Además, si
$K$ es el cuerpo de fracciones de $A$, el cuerpo de fracciones
$K'_0$ de $A'_0$ es radicial sobre $K$, puesto que el morfismo
$f$ es radicial. Como $A'_0$ es entero sobre $A$, su clausura
integral $B$ es también la clausura integral de $A$ en $K'_0$.
Pero, como $A$ es íntegramente cerrado, se sabe (Bourbaki,
\emph{Alg. comm.}, cap. V, § 2, n\textsuperscript{o} 3, lema 3) que $B$ es el
conjunto de los $x\in K'_0$ para los que alguna potencia
$p^m$-ésima (con $m$ suficientemente grande) pertenece a $A$
($p$ es el exponente característico de $K$). Además, no hay
más que un ideal primo de $B$ sobre cada ideal primo de $A$; en
particular, $B$ es un anillo local y su cuerpo residual es una
extensión radicial del de $A$ y, \emph{con mayor razón}, del de
$A'_0$. Esto prueba que $A'_0$ es geométricamente unirramoso
y, por consiguiente, también lo es $X'$.
\end{proof}
